%==============================================================================
% Home IoT Security — empirical study report
% Section skeleton mirroring the transportation IoT study
% ("On Security Vulnerabilities in Transportation IoT Devices").
% All section bodies are intentionally left blank.
%
% Build:  pdflatex home_iot_security_report.tex   (run twice for refs)
% Class:  IEEEtran (conference).
%==============================================================================
\documentclass[conference]{IEEEtran}
\IEEEoverridecommandlockouts

\usepackage{cite}
\usepackage{amsmath,amssymb,amsfonts}
\usepackage{graphicx}
\usepackage{booktabs}
\usepackage{array}
\usepackage{multirow}
\usepackage{xcolor}
\usepackage{url}
\usepackage[hidelinks]{hyperref}

% Inline placeholder for values the author must supply. Renders in bold red.
\newcommand{\stub}[1][?]{{\color{red}\textbf{[#1]}}}

\begin{document}

\title{On Security Vulnerabilities in Home IoT Devices}

\author{
\IEEEauthorblockN{Author Name}
\IEEEauthorblockA{Affiliation\\ City, State\\ email}
}

\maketitle

%==============================================================================
\begin{abstract}

\end{abstract}

\begin{IEEEkeywords}

\end{IEEEkeywords}

%==============================================================================
\section{Introduction}
\label{sec:intro}


%==============================================================================
\section{Related Work}
\label{sec:related}


%==============================================================================
\section{Our Methodology}
\label{sec:method}

In this section, we describe the methodology used in this paper. First, we
identified the device types belonging to each home IoT category and collected
the CVEs reported about them in NVD~\cite{ref-nvd}. We collected the CVEs using
two complementary search methods --- a vendor/brand search and a device-keyword
search --- and reconciled their results per category. Then, we extracted each
CVE's characteristics, including the corresponding CWEs and CVSS information.
Motivated by a previous study~\cite{ref-mcsoftware}, we grouped the individual
CWEs into the primary clusters of the CWE-888 Software Fault Patterns
view~\cite{ref-cwe888} in order to build a `big picture' of the security
vulnerabilities in home IoT devices. Finally, because keyword- and brand-based
NVD searches are text-based and therefore noisy, every candidate vulnerability
was subjected to a structured false-positive review before analysis.

The empirical work presented in this paper follows the case study (i.e.,
observational) approach based on publicly available data in NVD that are entered
and validated by experts and are widely used by researchers and practitioners.
Using an experimental approach (e.g., penetration testing) or a survey (e.g.,
collecting expert opinions) are complementary empirical approaches that are out
of the scope of this paper.

\subsection{Scope: what counts as a home IoT device}
\label{sec:method-scope}

We define a \emph{home IoT device} as an internet-connected sensor, appliance,
or embedded system deployed within a residential environment for the purpose of
monitoring, automation, or control, owned and maintained by non-expert consumers
without dedicated IT security oversight. For a device type to be in scope it must
satisfy all five definitional criteria: (1)~\emph{connectivity} --- it
communicates over a standard network protocol (TCP/IP, MQTT, CoAP, Zigbee, BLE);
(2)~\emph{device class} --- it is a special-purpose sensor, appliance, or
embedded system, not general-purpose IT (PC, phone, tablet, server, game
console); (3)~\emph{deployment context} --- it is intended for a private
residence, not primarily for enterprise or industrial use; (4)~\emph{function}
--- its primary purpose is to monitor, automate, or control the home environment
or its systems (climate, security, access, lighting, appliances, presence), with
media/entertainment, general computing, and communication explicitly excluded as
qualifying functions; and (5)~\emph{security context} --- it is owned and
maintained by consumers with no professional security administration. An
operational study-inclusion criterion further requires that the device have a
CPE-identifiable footprint in NVD and be subject to CVE disclosure. Because
connectivity (criterion~1) is satisfied by virtually every networked product, it
is the device's own function (criterion~4) and class (criterion~2) that
discriminate membership; a device that merely interoperates with home IoT is not
itself home IoT. Following criterion~4, entertainment devices such as game
consoles and streaming/smart TVs are excluded.

The study analyzes \stub[N] home IoT categories: \stub[final category list,
e.g.\ cameras, doorbells, baby monitors, door locks, alarms, thermostats, air
conditioners, fans, smart plugs, smart speakers, robot vacuums, refrigerators,
\ldots]. \stub[Note any categories that were defined but dropped or folded, e.g.\
sleep trackers (dominated by wearables) and the entertainment categories.]

\subsection{CVE collection from NVD}
\label{sec:method-collect}

CVEs are publicly known security vulnerabilities reported in NVD~\cite{ref-nvd}
and attributed to specific software applications or libraries. To maximize
recall, we collected CVEs related to home IoT devices using two complementary
discovery methods, each more reliable in a different respect:

\paragraph{Vendor/brand search} For each category we first compiled a list of
the manufacturers and brands that produce devices of that type, then searched NVD
for those vendor and brand names. This search was run offline against the local
NVD JSON year-feeds (\stub[year range, e.g.\ 2002--2026]), matching both the CVE
\emph{description text and the CPE (Common Platform Enumeration) strings}, so that
brand-in-CPE associations are captured. The brand lists were assembled from
\stub[brand-list sources, e.g.\ manufacturer directories / retailer catalogues /
targeted web searches per category]. Because brand names frequently overlap with
unrelated products, this method has higher recall but is more prone to false
positives.

\paragraph{Device-keyword search} In parallel, we searched NVD for generic
device-type keywords (e.g.\ ``ip camera'', ``baby monitor'', ``smart lock'') via
the NVD REST API v2.0. Note that the API's \texttt{keywordSearch} matches the CVE
\emph{description only} (it does not search CPE), so device-phrase keywords can
never match a brand encoded solely in a CPE string. Each CVE returned was
enriched with its CVSS score and severity, CWE identifier, and description.

\paragraph{Reconciliation} For each category we combined the two result sets
using set operations. The \emph{intersection} (vendor $\cap$ keyword) yields the
CVEs found by both methods. The \emph{difference} (vendor $-$ keyword) yields the
vendor CVEs the keyword search missed; this difference set is the primary input to
the false-positive review described in Section~\ref{sec:method-review}, both to
clean the dataset and to surface keywords the keyword search lacks.

We collected the data from NVD from \stub[start date] until \stub[end date]. In
total, we collected \stub[total CVEs] CVEs across the \stub[N] categories.
\stub[Note any vendors deliberately excluded for ambiguity --- analogous to the
transportation study's conservative exclusion of Qualcomm chips --- e.g.\ generic
chipsets or brand names that could not be confirmed as the device in question.]

\paragraph{Comparability caveat} The two searches are not perfectly comparable:
the keyword search hits the live NVD API (current data) and matches the
description only, whereas the vendor search runs over an offline year-feed
snapshot and matches description plus CPE. Consequently, some apparent ``gaps''
between the methods reflect data lag or the description-versus-CPE match surface
rather than genuine differences in coverage. \stub[State how this was handled for
the reported results --- e.g.\ whether both term lists were ultimately run
through the same engine against a single fixed NVD snapshot dated \stub[date], to
make the only difference the search terms.]

\subsection{CWE and CVSS data extraction}
\label{sec:method-extract}

Each CVE is typically attributed to one or more CWEs. However, some CVEs did not
have enough information to be associated with a CWE in NVD, or were not in NVD's
CWE mapping; these are listed as \texttt{NVD-CWE-noinfo} and
\texttt{NVD-CWE-Other}, respectively. Out of the total \stub[total CVEs] CVEs,
\stub[N] were \texttt{NVD-CWE-noinfo} and \stub[N] were \texttt{NVD-CWE-Other}.
These \stub[N] CVEs were included in our CVE-level analysis but did not contribute
to the CWE analysis. In total, \stub[total CWEs] CWEs were associated with the
\stub[total CVEs] CVEs; these were attributed to only \stub[unique CWEs] unique
CWEs.

For each CVE we also extracted the CVSS information. CVSS scores range from 0 to
10 and are computed from the vulnerability's exploitability and impact
metrics~\cite{ref-cvss}. CVSS severity is derived from the score: a score of 0
indicates the None severity level, and the Low, Medium, High, and Critical levels
correspond to score ranges of 0.1--3.9, 4.0--6.9, 7.0--8.9, and 9.0--10,
respectively. \stub[If both CVSS v2 and v3.x scores are present in the dataset,
state which version was used when both are available.] Lastly, for each CVE we
manually extracted information about the vulnerable components from the CVEs'
descriptions \stub[(confirm whether component extraction was performed, and how
components were named/normalized)].

\subsection{Mapping individual CWEs to the CWE-888 view}
\label{sec:method-cwe888}

CWE is a community-developed taxonomy of common weaknesses maintained by
MITRE~\cite{ref-cwe}. Because the CWE hierarchy is large and complex, several
views have been developed to organize it. We chose the CWE-888 Software Fault
Patterns view~\cite{ref-cwe888} to categorize the identified CWEs for two
reasons: (i)~it was used in~\cite{ref-mcsoftware}, which allows consistency and
direct comparison of results, and (ii)~it is focused on fault categorization, is
more technical, and captures a wider range of weaknesses than alternatives such
as the CWE-699 Software Development view~\cite{ref-cwe699}. We use the most recent
version of CWE-888 and focus only on its 23 primary clusters, which we refer to
as primary vulnerability classes.

If an identified CWE is itself part of the CWE-888 view, we mapped it directly to
its corresponding primary class. If it is not, we examined its parent CWEs and
mapped those to the CWE-888 view; a CWE whose parents fall under two different
primary classes therefore counts toward both. \stub[Report the over-counting
bookkeeping for this dataset: the number (and percentage) of CVEs whose CWE
mapped to two parents, and the number of CVEs that listed two distinct CWEs and
how those resolved across primary classes.] This approach builds a more complete
picture of the types of vulnerabilities present while limiting over-counting.

\subsection{False-positive classification}
\label{sec:method-review}

Because the vendor and keyword searches are purely text-based, generic brand names
and terms produce false positives (e.g.\ a brand string that also names an
unrelated commercial product, or marketing copy that merely mentions a device
category). Each candidate CVE in the difference set was therefore independently
classified as a true match for its category or a false positive.

To make this trustworthy, every row was judged \emph{independently and blind} by
\stub[five] reviewers --- two human researchers (\stub[Reviewer A] and
\stub[Reviewer B]) and three AI reviewers (Claude Code, ChatGPT Codex, and
Google's Gemini/Gemma). Blindness is enforced structurally: each reviewer works
on its own copy of the data containing only the raw fields plus its own empty
judgment columns, so no reviewer can see another's answer. All reviewers apply a
single shared rubric, judging \emph{only} from each row's \texttt{description} and
\texttt{cpe\_strings}.

Each reviewer assigns three fields per row: a \emph{Judgment} of \texttt{Yes}
(genuinely affects a home IoT device of this category), \texttt{No} (false
positive), or \texttt{Maybe} (genuinely ambiguous); a \emph{Confidence} of
\texttt{High} or \texttt{Low}; and a short, self-contained \emph{Reasoning}
required for \texttt{Low}-confidence and \texttt{Maybe} rows. A \texttt{Maybe} is
always paired with \texttt{Low} confidence. Low confidence is used when the device
is plausibly residential but primarily commercial/industrial, when the description
names the category but the CPE points at enterprise hardware, or when the CVE lies
in a software or protocol layer shared between home and non-home contexts. A
missing CPE string does not by itself force a \texttt{Maybe}: CPE data is
frequently absent on recent (2024--2026) CVEs due to NVD enrichment lag, so an
unambiguous description of a home device and a residential attack vector is
classified on its content.

The three AI reviewers' judgments are then merged. A row is flagged for human
adjudication when \stub[the human-review condition: both strong reviewers (Claude
and Codex) are Low confidence, or the three AI judgments are not unanimous].
Gemini is treated as a weaker third vote: its judgment counts toward unanimity,
but its self-reported confidence is recorded and excluded from the flag. A human
then adjudicates only the flagged rows, and each CVE's judgments are folded into a
single final label (\texttt{ai-consensus} for unflagged rows, \texttt{human} for
adjudicated rows). \stub[Report the outcome of the review --- e.g.\ the
false-positive rate per category, the number of rows flagged for human review, and
inter-reviewer agreement.]

%==============================================================================
\section{RQ1: Analysis of Vulnerability Classes in Home IoT Devices}
\label{sec:rq1}


\subsection{}
\label{sec:rq1-cat1}


\subsection{}
\label{sec:rq1-cat2}


\subsection{}
\label{sec:rq1-cat3}


\subsection{}
\label{sec:rq1-cat4}


%==============================================================================
\section{RQ2: Analysis of CVSS Scores and Severity}
\label{sec:rq2}


%==============================================================================
\section{Threats to Validity}
\label{sec:threats}


%==============================================================================
\section{Lessons Learned and Implications}
\label{sec:lessons}


%==============================================================================
\section{Conclusions}
\label{sec:conclusion}


%==============================================================================
\section*{Data Availability}


%==============================================================================
\begin{thebibliography}{00}
\bibitem{ref-nvd} NIST, ``National Vulnerability Database (NVD),''
  \url{https://nvd.nist.gov/}.
\bibitem{ref-mcsoftware} K. Goseva-Popstojanova and J. Tyo, ``Experience report:
  Security vulnerability profiles of mission critical software: Empirical
  analysis of security related bug reports,'' in \emph{IEEE 28th Int. Symp. on
  Software Reliability Engineering (ISSRE)}, 2017, pp.~152--163.
\bibitem{ref-cwe888} Common Weakness Enumeration, ``CWE-888: Software Fault
  Pattern (SFP) Clusters,''
  \url{https://cwe.mitre.org/data/definitions/888.html}.
\bibitem{ref-cwe699} Common Weakness Enumeration, ``CWE-699: Software
  Development,'' \url{https://cwe.mitre.org/data/definitions/699.html}.
\bibitem{ref-cwe} Common Weakness Enumeration, ``CWE,''
  \url{https://cwe.mitre.org/data/index.html}.
\bibitem{ref-cvss} FIRST, ``Common Vulnerability Scoring System (CVSS):
  Specification Document,'' \url{https://www.first.org/cvss/}.
\end{thebibliography}

\end{document}
